\chapter{Vectorization and Parallel Execution}\label{ch:vectorization-parallel-execution}

One instruction can do the work of many when data elements are independent.
A vector unit applies the same operation to a lane of values at once.
A parallel loop spreads iterations across threads or cores.
Both techniques exploit the same property: operations that do not interfere can share control.
This chapter shows how to test for that property and how to transform loops that have it.
We treat legality first, profitability second, and code shape last.
A fast loop that computes the wrong answer is not an optimization.

\section{What Independence Means}\label{sec:ch36-independence}

Two operations are independent when swapping their order changes nothing observable.
A loop is \term{vectorizable} when its iterations can run in lockstep lanes without changing results.
It is \term{parallelizable} when its iterations can run in any order, or concurrently, without changing results.
The obstacle is the \term{dependence}: a read or write whose value flows from one iteration to another.
Consider scaling an array by a constant.
Each iteration reads one element, multiplies, and writes it back.
No iteration touches another element, so all iterations are independent.
Now consider a running prefix sum where each iteration adds the previous total.
Iteration $i$ reads the value written by iteration $i - 1$.
That chain forbids both vector and parallel execution in its direct form.
Figure~\ref{fig:ch36-lanes} contrasts scalar execution with lane execution for the scaling example.
Scalar code handles one element per step.
Vector code handles four elements per step with one control decision.
Table~\ref{tab:ch36-patterns} lists common loop patterns with their dependence shape and verdict.
The table is qualitative because hardware widths and costs vary.
Element-wise maps usually succeed.
Reductions succeed with a combine step.
Recurrences with loop-carried flows fail without algorithmic change.

\begin{figure}[htbp]
\centering
\begin{tikzpicture}[node distance=0.9cm and 1.2cm]
\node[basic block] (s1) {scalar\\ one element};
\node[basic block, below=of s1] (s2) {scalar\\ next element};
\node[basic block, right=of s1, minimum width=2.6cm] (v1) {vector\\ lanes 0--3};
\node[basic block, below=of v1, minimum width=2.6cm] (v2) {vector\\ lanes 4--7};
\draw[cfg edge] (s1) -- (s2);
\draw[cfg edge] (v1) -- (v2);
\end{tikzpicture}
\caption{Scalar code steps through elements one at a time while vector code covers four lanes per step, which pays only when lanes are independent.}
\label{fig:ch36-lanes}
\end{figure}

\begin{table}[htbp]
\centering
\caption{Loop patterns differ in dependence shape and in their threading verdict.}
\label{tab:ch36-patterns}
\begin{tabular}{lll}
\toprule
Pattern & Dependence & Verdict \\
\midrule
Element-wise map & None & Vectorize and parallelize \\
Reduction & Combine only & Tree reduce then combine \\
Stencil, shift 1 & Bounded neighbors & Vectorize with alignment plan \\
Recurrence & Loop-carried chain & Reject without rewrite \\
\bottomrule
\end{tabular}
\end{table}

\section{Testing Legality: Distance and Direction}\label{sec:ch36-legality}

A \term{dependence distance} measures how many iterations separate a write from its later read.
Distance zero means the flow stays inside one iteration and rarely blocks transformation.
Positive distances mean later iterations consume earlier results.
A \term{direction vector} records per-loop whether the flow goes forward, backward, or stays put.
Algorithm~\ref{alg:ch36-vector-check} gives the legality test for simple affine loops over arrays.
It assumes subscripts are linear functions of loop indices with known coefficients.
It requires alias facts for pointer accesses: two accesses that may overlap are treated as a dependence.
It returns a verdict of vectorize, parallelize, or reject with the blocking dependence.
The test walks each pair of accesses to the same array.
It solves for index pairs that touch the same element.
It converts the iteration gap into a distance and a direction.
Zero distances everywhere permit vectorization at the innermost level.
Parallelization needs the stronger property that no positive distance crosses thread boundaries, possibly after privatization.
Privatization gives each thread or lane a private copy of a temporary that would otherwise serialize the loop.
Reductions receive special handling: the combine operator must be associative, with a documented identity and a tree that preserves the sequential result up to stated rounding.
Floating-point reassociation changes results, so the compiler needs explicit permission before reordering those operations.

\begin{algorithm}[htbp]
\caption{Legality check for vectorizing an affine loop nest}\label{alg:ch36-vector-check}
\begin{pseudocode}
\Require affine loop nest $L$, complete scalar and memory dependence information, operations safe for proposed lane execution
\Ensure verdict: vectorize, parallelize, or reject with cause
\State $D \gets$ all dependences with distance and direction
\If{any pair may alias without disambiguation}
\State \Return reject with cause
\EndIf
\If{there is a reduction and all loop-carried dependences belong to it, with a legal combine order}
\State \Return parallelize with tree reduction
\EndIf
\If{any remaining loop-carried dependence exists on the target loop}
\State \Return reject with cause
\EndIf
\State \Return vectorize innermost loop
\end{pseudocode}
\end{algorithm}

\section{Worked Example: Distances in Three Loops}\label{sec:ch36-worked}

We compute distances for three loops over arrays of length $8$ with indices $0$ through $7$.
Loop one scales in place: \code{a[i] = a[i] + 1}.
The read and write in iteration $i$ touch element $i$.
Equal indices give distance $0$.
All distances are zero, so Algorithm~\ref{alg:ch36-vector-check} returns vectorize.
With width $4$, iterations $0$--$3$ form one vector and $4$--$7$ form the next.
Loop two shifts right: \code{a[i] = a[i - 1] + 1} for $i$ from $1$ to $7$.
Iteration $i$ reads the element written by iteration $i - 1$.
The distance is $1$ in the forward direction.
The algorithm returns reject for direct vectorization because a positive distance crosses lanes.
Loop three accumulates a total: \code{s = s + a[i]}.
Every iteration reads and writes $s$.
The distance is $1$ on a scalar, not an array element.
The algorithm returns parallelize with tree reduction when addition is the documented combine.
Hand arithmetic confirms the reduction shape: partial sums over $0$--$3$ and $4$--$7$ combine once at the end.
The listing below shows the vector shape for loop one.
\begin{irlisting}
; vector scale, width 4, length 8
; v0 = load a[0..3]
; v1 = add v0, 1
; store v1 -> a[0..3]
; v2 = load a[4..7]
; v3 = add v2, 1
; store v3 -> a[4..7]
\end{irlisting}
Remainder handling follows the same rule.
A length of $10$ uses two full vectors and a scalar tail of $2$.
The tail executes the original scalar body with indices $8$ and $9$.

\section{Profitability, Alignment, and Tail Code}\label{sec:ch36-profit}

Legality permits transformation but does not require it.
\term{Profitability} weighs vector speedup against masks, tails, gathers, and extra memory traffic.
Short trip counts may not fill one vector.
Conditional bodies need masks that idle some lanes.
Noncontiguous access needs gathers or a layout change.
A simple model guides the decision: gain comes from lanes per instruction, cost comes from setup plus tail plus memory stalls.
\begin{ccnote}[Measure the memory]
Vector arithmetic is cheap and memory is not.
A loop that misses the cache at every step stays slow in any lane width.
\end{ccnote}
Alignment adds a concrete constraint.
Vector loads run fastest at aligned addresses and fault or slow down otherwise.
The compiler peels a few iterations to reach alignment, aligns the main body, and handles the tail separately.
Alias uncertainty adds another.
When two pointers may overlap, the compiler versions the loop: a fast vector body for the disjoint case and a scalar body otherwise, selected by a run-time check.
The \enquote{version then vectorize} pattern keeps the common case fast without assuming facts the analysis cannot prove.

\section{Tiling for Caches Before Parallelism}\label{sec:ch36-tiling}

Threads do not fix memory bottlenecks.
A loop that streams past the cache on one thread streams past the cache on every thread, only with more contention.
\term{Loop tiling} restructures nested loops so each block of iterations reuses data while it is still cached.
The classic case is matrix multiplication over $N$-by-$N$ inputs.
The triple nest performs $O(N^3)$ arithmetic operations.
Its cache misses depend on loop order, layout, cache capacity, and cache-line size; operation count alone does not determine them.
Tiling seeks to reuse resident submatrices before eviction, and any quantitative transfer bound needs an explicit memory model.
The block size $B$ comes from cache capacity, line size, and associativity, not from the algorithm.
Tiling composes with vectorization and threading in a fixed order.
Tile first so working sets fit private caches.
Vectorize the innermost body so each thread issues wide operations.
Parallelize the outer tile loop so threads share no working set.
Consider an $8$-by-$8$ product with $4$-by-$4$ tiles.
Four $4$-by-$4$ tiles cover the output.
Each output tile traverses two reduction-dimension tiles, reading a pair of $4$-by-$4$ input blocks and performing $4^3=64$ multiply-add updates for each pair.
Thus each output tile performs $128$ updates and the four tiles perform $512$, equal to $8^3$.
Counting multiplication and addition separately gives $1024$ scalar arithmetic operations; fusion into machine instructions is a separate target decision.
The bounded integer model in \code{example-code/refinement\_checks.c} compares this $8$-by-$8$ tiled product with a naive nest and checks all $512$ updates.
A general implementation should also test prime-dimension inputs that leave partial edge tiles.
The \enquote{edges match} check catches the defect where edge tiles read past the input or skip the final rows.

\section{Implementation Notes and Complexity}\label{sec:ch36-implementation}

For $P$ candidate access pairs, dependence testing costs the sum of $P$ query costs; affine integer feasibility is not a constant-time operation in general.
Exact testing is expensive in general, so compilers use layered tests: fast sufficient checks first, exact solving only for survivors.
Vector code generation is linear in the loop body size plus tail and peel code.
Threading adds synchronization: barriers cost $O(T)$ wakeups for $T$ threads and reductions cost $O(\log T)$ combine depth in the ideal case.
Memory bandwidth caps both.
Production compilers vectorize innermost loops, parallelize outer loops, and tile for caches before either.
They respect floating-point ordering rules and require explicit flags before reassociation.
They keep scalar, masked, and threaded versions behind run-time checks when static facts run out.
A common pitfall vectorizes a loop with a hidden carried dependence through a scalar temporary.
Privatization must cover every such temporary or the lanes interfere.
Another pitfall parallelizes a loop with an unprotected reduction and reports nondeterministic sums.
The tree reduction with a stated combine order prevents that surprise.
Test each transformed loop against its scalar result on full, empty, and prime-length inputs.
% TODO(cite): need a primary reference for direction-vector legality proofs.

\section{Extend profitability to placement and heterogeneous targets}
\label{sec:ch36-placement-link}

An independent loop may have several legal implementations: scalar, vector, threaded, or device offloaded.
Choosing among them requires a cost model that includes setup, transfers, synchronization, and data residency as well as arithmetic throughput.
For synthetic costs $4n$ on a CPU and $120+n$ on a device, the device wins only above $n=40$, despite its lower per-element coefficient.
The legality proof does not decide this threshold.

Chapter~\ref{ch:profile-spatial-compilation} works that comparison alongside learned cost fitting, uncertain speculation benefits, and a small exact placement problem.
It also distinguishes temporal affinity from dependence: operations often observed close together may be worth placing nearby, but observation alone does not permit exchanging their execution order.

\section{Exercises}\label{sec:ch36-exercises}

\begin{enumerate}
\item Compute distances for \code{a[i] = a[i] + a[i - 2]} over $i$ from $2$ to $7$, state the direction, and give the verdict of Algorithm~\ref{alg:ch36-vector-check}.\label{ex:ch36-distance-math}
\item Plan aligned vector code for length $10$ at width $4$ starting at a misaligned address, giving peel count, main-vector count, and tail count.\label{ex:ch36-tail-plan}
\item Write a reduction test for a sum over $1000$ values that detects a missing combine step, describing the thread split and the expected partial sums.\label{ex:ch36-reduction-test}
\end{enumerate}

\section{Further Reading}\label{sec:ch36-reading}

Appel covers loop structure and dataflow facts that dependence testing builds upon~\cite{appel-1998-modern}.
Morgan relates loop shape and memory behavior to optimization choices~\cite{morgan-1998-optimizing}.
Safonov discusses preserving program meaning under aggressive transformation~\cite{safonov-2010-trustworthy}.
The C standard draft defines the aliasing and ordering rules that vectorization must respect~\cite{wg14-2011-c11}.
